% ---------------------------------------------------------------------------
% Pseudo-satellite assimilation in the CHIRPS OSSE: a defect, not a result.
%
% DRAFTING NOTE: recorded in full, as with docs/ablation_v3.tex, so the section
% can be cut down later without re-deriving anything. Source data:
%   data/processed/osse_paper/gauges_realistic_40/
%   data/processed/osse_paper/satellite_realistic_40/
%   data/processed/osse_paper/simultaneous_realistic_40/
%   data/processed/osse_paper/paper/  (scripts/24_osse_paper_suite.py output)
% Consequence: slurm/osse_gauges_bmd.sbatch runs gauges only until this is
% resolved.
% ---------------------------------------------------------------------------

\section{Pseudo-satellite assimilation degrades the analysis: diagnosis}
\label{sec:pseudo-satellite-defect}

\subsection{What was run}

The CHIRPS OSSE over 2021--2024 supplies a $0.05^\circ$ nature run, pseudo-gauge
values at real BMD coordinates, and exact nested $0.1^\circ$ block means as
pseudo-IMERG. Three matched arms share dates, seeds and holdout: gauges only,
pseudo-IMERG only, and simultaneous. Scoring is at withheld pseudo-gauges,
with gridded diagnostics at the full field, at the $0.1^\circ$ footprints, and
on the subgrid residual left after removing each field's own footprint mean.

Both observation types are drawn from the nature run and then perturbed. Under
\texttt{--obs-error realistic}, which is what the arms below use, gauges receive
\emph{white} noise of standard deviation
$\sqrt{\sigma_{\mathrm{obs}}^2 + \mathrm{rep}^2} = \sqrt{0.10^2 + 0.25^2} =
0.269$ in transformed units, and the pseudo-satellite receives a
\emph{spatially correlated} error described in
Section~\ref{sec:pseudo-satellite-defect} below. Under
\texttt{--obs-error exact} neither is perturbed at all: observed values are the
CHIRPS pseudo-data unchanged, and $R = 0.05$ is a numerical regulariser, since
the likelihood variance $V(t) = R + \gamma(1-t)^2/t^2$ is singular as
$t \to 1$.

The distinction matters for reading the table. These are \emph{not} perfect
observations, and the gauge arm's $+14\%$ is skill obtained despite a
realistic error, not an upper bound.

\subsection{The numbers}

\begin{table}[t]
\centering
\caption{OSSE arms against their own matched background. Percentages are
improvements, so negative means the analysis is WORSE than the prior it
started from. $\Delta r$ is analysis minus background correlation.}
\label{tab:pseudo-satellite}
\small
\begin{tabular}{lrrr}
\hline
 & gauges only & pseudo-IMERG only & simultaneous \\
\hline
Withheld CRPSS (\%)        & $+14.03$ & $-68.42$  & $-38.65$ \\
Withheld RMSE gain (\%)    & $+10.05$ & $-146.03$ & $-91.04$ \\
Field CRPSS (\%)           & $+25.96$ & $-34.17$  & $-2.11$  \\
$0.1^\circ$ CRPSS (\%)     & $+27.59$ & $-36.33$  & $-2.75$  \\
Subgrid CRPSS (\%)         & $+3.65$  & $-5.41$   & $-1.19$  \\
Subgrid RMSE gain (\%)     & $+0.53$  & $-5.31$   & $-2.84$  \\
\hline
$\Delta r$, withheld gauges & $+0.110$ & $-0.058$ & $-0.049$ \\
$\Delta r$, full field      & $+0.130$ & $-0.017$ & $+0.040$ \\
$\Delta r$, footprints      & $+0.133$ & $-0.019$ & $+0.040$ \\
$\Delta r$, subgrid         & $+0.018$ & $+0.038$ & $+0.047$ \\
\hline
Withheld coverage$_{90}$    & 0.581 & 0.688 & 0.667 \\
Withheld spread/skill       & 0.698 & 0.595 & 0.643 \\
\hline
CRPS at withheld gauges (mm\,day$^{-1}$) & & & \\
\quad background            & 4.22 & 4.22 & 4.22 \\
\quad analysis              & \textbf{3.64} & 7.14 & 5.86 \\
\hline
\end{tabular}
\end{table}

\subsection{Why this is a defect and not a result}

\textbf{The gauge arm works.} It improves CRPS by $14\%$ at withheld
pseudo-gauges, the full field by $26\%$, the footprint scale by $28\%$, and
raises correlation by $+0.11$ at withheld stations. The sampler, the guidance,
the holdout and the scoring are therefore all functioning; whatever is wrong is
specific to the pseudo-satellite path.

\textbf{An uninformative observation cannot do this.} Assimilating an
observation that carries no usable information should drive the analysis toward
the prior, converging on zero improvement. It cannot push the analysis
\emph{below} the prior by $68\%$ in CRPS and $146\%$ in RMSE. Degradation of
that magnitude requires the likelihood to be actively pulling the state toward
somewhere the truth is not. Since the pseudo-observations \emph{are} block means
of the truth, the error must lie in how those values are generated, transformed,
or compared against the operator's prediction --- not in the observation's
information content.

\subsection{The signature points at magnitude, not noise}

The discriminating comparison is between the correlation changes and the
RMSE changes, because correlation is invariant to a constant offset and to a
uniform scaling while RMSE is not.

Correlation is \emph{approximately unchanged} where the observation acts: at
the footprint scale the pseudo-IMERG arm moves $r$ by only $-0.019$, and at the
full field by $-0.017$. Over the same fields RMSE degrades by $81\%$ and CRPS by
$34$--$36\%$. A field can only keep its correlation while losing that much RMSE
if its \emph{values} are wrong and its \emph{pattern} is broadly right.

Random observation noise does not behave this way. Noise degrades pattern and
amplitude together, so $r$ would fall alongside RMSE. What the table shows
instead is the signature of a systematic offset or scale error applied
consistently across footprints.

Two further details are consistent with that reading and with nothing else
obvious. First, the subgrid residual --- the component the satellite does
\emph{not} constrain --- is the one place pseudo-IMERG \emph{improves}
correlation ($+0.038$, against $+0.018$ for gauges alone), which is what one
expects if a large systematic increment at footprint scale leaves the
within-footprint structure comparatively untouched. Second, coverage$_{90}$
\emph{rises} to $0.688$ from the gauge arm's $0.581$ while spread/skill
\emph{falls} to $0.595$: the ensemble is not becoming better calibrated, its
error is growing faster than its spread.

\paragraph{A caveat on the correlation argument.}
At the withheld gauges themselves $\Delta r$ is $-0.058$, a larger drop than at
field or footprint scale. A pure offset would leave point correlation untouched,
so the mechanism is not \emph{only} an offset. An intensity-\emph{dependent}
distortion --- one whose size grows with rainfall amount --- would produce
exactly this: near-preserved correlation at aggregate scales, with a measurable
loss at points where intensities vary most. That is the form the leading
hypothesis below takes.

\subsection{Cause, established from the code}

An earlier draft of this section proposed that the pseudo-observations were
built by averaging \emph{transformed} CHIRPS while the operator averages in
physical space, making $T(\overline{x})$ and $\overline{T(x)}$ inconsistent.
\textbf{That hypothesis is wrong} and is recorded here only so it is not
proposed again. \texttt{scripts/10\_osse.py} computes the pseudo-satellite truth
as \texttt{block\_mean\_mm(truth, \ldots)} followed by
\texttt{transform.forward(\ldots)} --- a physical-space mean, then the transform
--- which is exactly what \texttt{PhysicalBlockAverageObsOperator} predicts. The
two are consistent.

The actual cause is that \textbf{the correlation inflation applied to $R$ is
also used to manufacture the observation error}. In sequence:

\begin{enumerate}
  \item $R$ for the satellite block is built from the configured
    $\sigma_{\mathrm{obs}} = 0.35$ and representativeness $0.10$, giving a
    standard deviation of $\sqrt{0.35^2 + 0.10^2} = 0.364$ in transformed units.
  \item It is then inflated in place by
    $\max\!\left(1,\; 2\pi(\ell/s)^2\right)$ with correlation length
    $\ell = 3$ cells and stride $s = 3$, i.e.\ by $6.28\times$ in variance:
\begin{verbatim}
if use_satellite and satellite_r_inflation != 1.0:
    R[len(active_assim_idx):] *= satellite_r_inflation
\end{verbatim}
  \item The pseudo-IMERG product is then \emph{generated} by drawing from that
    same inflated $R$, with spatial correlation imposed:
\begin{verbatim}
satellite_observed = perturb_observations(
    satellite_truth_transformed, satellite_R, 1, corr_blocks=[...])
\end{verbatim}
\end{enumerate}

The synthetic satellite is therefore corrupted with a \emph{spatially
correlated} error of standard deviation $0.364 \times \sqrt{6.28} = 0.91$
transformed units, against the gauges' \emph{white} $0.269$. The pseudo-IMERG
product carries roughly $3.4\times$ the gauge error, and because that error is
correlated over $\sim\!30$~km it does not average away across footprints.

This is a design error, not a tuning choice. The correlation inflation exists to
\emph{down-weight} observations whose errors are not independent. Using it to
build the observation means that asking for a more conservative $R$ makes the
synthetic product worse, and the two effects are then counted twice: once in
the data and once in the weight. An OSSE should generate the observation with
its \emph{true} error and assimilate it under whatever $R$ is being tested.

\paragraph{This explains every feature of Table~\ref{tab:pseudo-satellite}.}
A large correlated error displaces whole regions coherently, which destroys
RMSE while leaving the broad pattern --- and hence $r$ --- largely intact. And
it explains the one anomaly a magnitude-error story could not: subtracting each
field's own $0.1^\circ$ footprint mean removes most of a footprint-scale
correlated error, so the \emph{subgrid} residual is comparatively clean and its
correlation \emph{improves} ($+0.038$), while everything at and above footprint
scale collapses.

\paragraph{The same mechanism, now seen twice.}
This is the $\sqrt{R}$ pathology already established on real data, where stride
1 at $0.1^\circ$ degraded CRPS by $+2.388~[+1.824, +2.956]$ because
\texttt{perturb\_observations} draws with $\mathrm{sd} = \sqrt{R}$ and a
$56.5\times$ inflation therefore multiplied the perturbation amplitude by
$7.5$. It has now appeared in two independent code paths --- once when
perturbing ensemble members, once when constructing a synthetic observation ---
which makes it a property of how $R$ is used throughout this system rather than
a local slip.

\subsection{The fix, and how to confirm it}

Generate the pseudo-observation from the \emph{uninflated} error and keep the
inflation for the assimilation weight only. Concretely, pass the pre-inflation
$R$ to \texttt{perturb\_observations} at the product-generation step while
leaving the inflated $R$ in the likelihood.

The confirmation is cheap and needs no GPU: re-run the satellite arm with
\texttt{--obs-error exact}, which bypasses \texttt{perturb\_observations}
entirely (\texttt{satellite\_observed = satellite\_truth\_transformed.copy()})
and is already guarded by an invariant check against the same-day CHIRPS block
mean. If the exact satellite arm improves on its background while the realistic
one does not, the observation-generation path is confirmed as the cause and the
operator, geometry and likelihood are all exonerated.

\subsection{Consequence for the paper}

The pseudo-satellite contaminates every arm it enters: the simultaneous arm is
also below its background at withheld gauges ($-38.65\%$), so no arm containing
pseudo-IMERG can currently be reported. The OSSE has therefore been re-run
gauges-only on the real BMD geometry
(\texttt{slurm/osse\_gauges\_bmd.sbatch}), where the machinery demonstrably
works, and satellite arms are excluded until the diagnosis above is resolved.

This does \emph{not} affect the real-data results in
Section~\ref{sec:ingestion}, which use actual IMERG through a different code
path (\texttt{scripts/15\_bmd\_month\_example.py}) and never touch the OSSE's
pseudo-observation generator. It does, however, mean the OSSE cannot yet be
used to explain the real-data null --- which was its main purpose in the
argument --- and that experiment is deferred rather than abandoned.
